\documentclass[aspectratio=169,11pt]{beamer}
\usepackage[utf8]{inputenc}
\usepackage[T1]{fontenc}
\usepackage[spanish]{babel}
\usepackage{lmodern}
\usepackage{booktabs}
\usepackage{tabularx}
\usepackage{hyperref}
\usepackage{enumitem}

\usetheme{Madrid}
\usecolortheme{default}
\setbeamertemplate{navigation symbols}{}
\setbeamertemplate{footline}[frame number]

\title{IA agéntica para la creación de contenido académico}
\subtitle{Sesión 2 --- Briefs, prompts de sistema y orquestación}
\author{Juliho Castillo}
\institute{CADi 40\,h · Escuela de Ingeniería y Ciencias\\
Tecnológico de Monterrey}
\date{Clave sugerida: \texttt{CADI-EIC-IA-AGENTICA-01}}

\begin{document}

% --- 1 Portada ---
\begin{frame}
\titlepage
\end{frame}

% --- 2 Agenda ---
\begin{frame}{Agenda (3\,h) y producto E2}
\begin{tabularx}{\textwidth}{@{}lclX@{}}
\toprule
Bloque & Tiempo & Contenido \\
\midrule
Recap & 10 min & Integridad y temas E1 \\
Brief & 30 min & Contrato pedagógico \\
Prompt y multi-rol & 25 min & Anatomía y roles \\
Demo orquestación & 25 min & HITL y criterio de parada \\
Taller & 50 min & Brief + prompt + diagrama \\
Cierre & 20 min & Puesta en común; TI \\
\bottomrule
\end{tabularx}
\vspace{0.5em}
\textbf{E2:} brief + prompt de sistema v0 + diagrama de orquestación.
\end{frame}

% --- 3 RA ---
\begin{frame}{Resultados de aprendizaje (sesión)}
Al finalizar, usted podrá:
\begin{enumerate}
\item Completar un \textbf{brief} de contenido específico para material EIC
\item Diseñar un \textbf{prompt de sistema} (rol, restricciones, formato, parada)
\item Proponer \textbf{orquestación} de 2--4 roles con aprobaciones humanas
\item Incluir instrucciones de \textbf{no fabricar citas} y señalar incertidumbre
\end{enumerate}
\end{frame}

% --- 4 Puente E1 ---
\begin{frame}{Puente desde E1}
\begin{itemize}
\item Trajo: mapa + nota de integridad + tema EIC tentativo
\item Hoy: convertir esa intención en \textbf{contrato} (brief) y \textbf{máquina} (prompt + flujo)
\item El brief alimenta al planificador/redactor; el diagrama fija los HITL
\end{itemize}
\begin{block}{Recordatorio}
Política de integridad / uso de IA $\rightarrow$ verificar con \textbf{CEDDIE} \emph{(enlace y vigencia por edición)}.
\end{block}
\end{frame}

% --- 5 Brief contrato ---
\begin{frame}{El brief como contrato pedagógico}
\begin{itemize}
\item Contrato entre \textbf{docente} y \textbf{sistema} (no un ``prompt largo'')
\item Fija audiencia, alcance, tono, restricciones y criterios de éxito
\item Reduce ``ruido elegante'' y regeneraciones inútiles
\item Marque explícitamente: ``El humano decide siempre sobre X''
\end{itemize}
\begin{alertblock}{Mensaje clave}
Sin brief sólido, la orquestación solo multiplica el caos.
\end{alertblock}
\end{frame}

% --- 6 Campos brief ---
\begin{frame}{Campos mínimos de un buen brief EIC}
\begin{enumerate}
\item Identificación (artefacto, curso, idioma)
\item Audiencia y conocimientos previos
\item Objetivos / propósito (verbo + contenido)
\item Alcance: incluye / excluye / extensión
\item Tono y estilo; notación preferida
\item Restricciones para el agente + criterios de éxito
\end{enumerate}
Plantilla: \texttt{tex/plantillas/brief-contenido.tex}
\end{frame}

% --- 7 Prompt anatomía ---
\begin{frame}{Anatomía del prompt de sistema}
\begin{itemize}
\item \textbf{Rol:} planificador, redactor, revisor, verificador\ldots
\item \textbf{Misión:} qué producir y para quién
\item \textbf{Reglas duras:} no inventar citas/datos; \texttt{[VERIFICAR]}
\item \textbf{Formato de salida:} español (México); supuestos; pendientes
\item \textbf{Criterio de parada:} cuándo devolver control al humano
\end{itemize}
Plantilla: \texttt{tex/plantillas/prompt-sistema-agente.tex}
\end{frame}

% --- 8 Restricciones STEM ---
\begin{frame}{Restricciones STEM (obligatorias)}
Incorpore en el prompt:
\begin{itemize}
\item No fabricar citas, DOI, estadísticas ni URLs
\item Señalar incertidumbre con \texttt{[VERIFICAR: \ldots]}
\item Respetar notación, unidades y glosario del curso
\item No solicitar ni pegar datos personales sensibles
\item Distinguir hecho, interpretación y opinión
\end{itemize}
\begin{alertblock}{STEM académico}
Estas reglas son el mínimo; no son ``opcionales''.
\end{alertblock}
\end{frame}

% --- 9 Patrones orquestación ---
\begin{frame}{Patrones de orquestación}
\begin{itemize}
\item \textbf{Secuencial:} plan $\rightarrow$ borrador $\rightarrow$ revisión
\item \textbf{Crítico-humano:} el humano es un nodo obligatorio
\item \textbf{Paralelo ligero:} dos roles $\rightarrow$ fusión supervisada
\end{itemize}
\begin{center}
\texttt{Usuario $\rightarrow$ Orquestador $\rightarrow$ Roles $\rightarrow$ Herramientas $\rightarrow$ Humano}
\end{center}
Orquestación $>$ prompt único para obras largas o multi-artefacto.
\end{frame}

% --- 10 Mapa roles ---
\begin{frame}{Mapa de roles académicos (2--4)}
\begin{tabularx}{\textwidth}{@{}lX@{}}
\toprule
Rol & Entrega típica \\
\midrule
Planificador & Outline / arquitectura alineada a RA \\
Redactor & Sección o bloque con restricciones \\
Revisor & Hallazgos priorizados (no reescritura ciega) \\
Verificador & Notación, unidades, citas, consistencia \\
\bottomrule
\end{tabularx}
\vspace{0.8em}
Empiece con 2--3 roles; agregue verificador cuando haya STEM denso.
\end{frame}

% --- 11 Parada y handoff ---
\begin{frame}{Criterio de parada y handoff}
El agente se detiene y pide humano cuando:
\begin{itemize}
\item Completó el encargo del brief
\item Necesita una decisión pedagógica o disciplinar
\item Detecta riesgo de fabricar información
\item Hay ambigüedad: ofrece 2 opciones etiquetadas
\end{itemize}
\begin{block}{Handoff limpio}
Salida estructurada + supuestos + pendientes de verificación.
\end{block}
\end{frame}

% --- 12 Demo ---
\begin{frame}{Demo --- orquestación con HITL visible}
Flujo en vivo (ejemplo STEM neutro EIC):
\begin{enumerate}
\item Brief corto $\rightarrow$ planificador propone outline
\item \textbf{HITL\,1:} humano aprueba / corta nodos
\item Redactor produce un bloque con restricciones
\item \textbf{HITL\,2:} humano marca \texttt{[VERIFICAR]} y edita
\end{enumerate}
Observe: restricciones + parada + decisiones documentadas.
\end{frame}

% --- 13 Anti-patrones ---
\begin{frame}{Anti-patrones de brief y orquestación}
\begin{itemize}
\item Brief genérico (``haz un capítulo bueno'')
\item Un solo mega-prompt sin roles ni parada
\item Cero restricciones STEM
\item Aceptar la primera salida multi-rol sin fusión humana
\item Olvidar el diagrama: nadie sabe dónde está el HITL
\end{itemize}
\end{frame}

% --- 14 Actividad ---
\begin{frame}{Actividad de taller ($\sim$50 min)}
\textbf{Entregable E2} (un solo archivo):
\begin{enumerate}
\item \textbf{Parte A --- Brief} (20 min): campos mínimos + ``humano decide X''
\item \textbf{Parte B --- Prompt} (15 min): rol, misión, reglas, formato, parada
\item \textbf{Parte C --- Orquestación} (15 min): 2--4 roles, flujo, $\geq$2 HITL
\end{enumerate}
Nombre sugerido: \texttt{Apellido\_Sesion2\_brief-prompt-orquestacion}
\end{frame}

% --- 15 Plantillas ---
\begin{frame}{Plantillas a usar hoy}
\begin{itemize}
\item \texttt{tex/plantillas/brief-contenido.tex}
\item \texttt{tex/plantillas/prompt-sistema-agente.tex}
\item Su mapa E1 como base del diagrama
\end{itemize}
\begin{block}{Prioridad del taller}
Calidad del brief y del diagrama $>$ volumen de texto generado.
\end{block}
\end{frame}

% --- 16 Rúbrica ---
\begin{frame}{Rúbrica rápida (E2)}
\begin{tabularx}{\textwidth}{@{}lXX@{}}
\toprule
Criterio & Bien & A mejorar \\
\midrule
Alineación EIC & Audiencia / notación / objetivos & Genérico \\
HITL & $\geq$2 puntos y \texttt{[VERIFICAR]} & Aceptación pasiva \\
Reutilización & Sirve a más de una unidad & Solo ad hoc \\
Integridad & No-citas + parada explícita & Omite restricciones \\
\bottomrule
\end{tabularx}
\end{frame}

% --- 17 Takeaways ---
\begin{frame}{Takeaways de la sesión}
\begin{enumerate}
\item \textbf{El brief es el contrato} entre docente y sistema
\item \textbf{Orquestación $>$ prompt único} en obras largas
\item \textbf{Criterio de parada} deja claro cuándo entra el humano
\item Restricciones STEM reducen alucinaciones ``elegantes''
\end{enumerate}
\end{frame}

% --- 18 TI ---
\begin{frame}{Trabajo independiente --- bloque B}
Antes de la sesión 3 ($\sim$parte del presupuesto de 16\,h):
\begin{itemize}
\item Practicar plantillas de brief y prompt de sistema
\item Refinar E2 con su tema EIC real
\item Traer idea de \textbf{arquitectura de obra} (capítulos / unidades)
\end{itemize}
Canal de dudas: \emph{(definir por edición)}.
\end{frame}

% --- 19 Prep sesión 3 ---
\begin{frame}{Prep sesión 3 --- Libros/ebooks I}
Próxima sesión (3\,h):
\begin{itemize}
\item Arquitectura pedagógica de una obra técnica
\item Outline 6--10 nodos alineado a RA
\item Mapa preliminar de notación / glosario
\end{itemize}
Traiga brief + diagramas E2 y su tema EIC.
\end{frame}

% --- 20 Cierre ---
\begin{frame}{Gracias}
\begin{center}
{\Large Juliho Castillo}\\[0.5em]
Escuela de Ingeniería y Ciencias\\[1em]
CADi · IA agéntica para la creación de contenido académico\\[1.5em]
{\small \emph{Verificar políticas de integridad y uso de IA con CEDDIE antes de cada edición.}}
\end{center}
\end{frame}

\end{document}
